% !TeX TXS-program:compile = txs:///pdflatex/[-shell-escape]
% !TeX TXS-program:pdflatex = pdflatex -synctex=1 -interaction=nonstopmode --shell-escape %.tex
\section{The Stack, Revisited}\label{sec:10the_stack_revisited}

\sectiontoc

\subsection{Reaching Into the Caller's Frame}\label{subsec:10-01reaching_into_the_callers_frame}

La pila almacena las variables locales de las funciones: datos que estas utilizan para funciones que no son
necesariamente necesarias para otras funciones del programa.

Por lo tanto, tú podrías acceder al marco de una función para poder conocer
los valores de sus variables, y la dirección de retorno.

\begingroup
\setlength{\dashlinedash}{1pt}
\setlength{\dashlinegap}{2pt}

\begin{tabular}{:c:p{8cm}}
	\multicolumn{1}{c}{\ }        & \makecell{[smaller addresses] \\ \( \leftarrow \)rsp, immediately before `call solve` }\\
	\cdashline{1-1}
	\makecell{caller's local region                               \\ \dots your flag is in here \dots } &                                                                   \\
	\cdashline{1-1}
	caller's saved rbp            &                               \\
	\cdashline{1-1}
	return address (back to main) &                               \\
	\cdashline{1-1}
	\dots main's frame \dots      &                               \\
	\cdashline{1-1}
	\multicolumn{1}{c}{\ }        & [large addresses]             \\
\end{tabular}
\endgroup

La instrucción \mintinline{ASM}{call solve} realiza 2 acciones:

\begin{enumerate}
	\item Empula la dirección de retorno a la pila (8bytes), por lo que la dirección de retorno acaba ocupando una dirección menor que la que ya había en la pantalla.
	\item Salta a tu código.
\end{enumerate}

El primer paso es fundamental, puesto que la pila crece al revéz. La instrucción
\mintinline{ASM}{pop} suma 8 a \mintinline{ASM}{rsp} y la instrucción
\mintinline{ASM}{push} resta 8, equivalencias:

\begin{itemize}
	\item \mintinline{ASM}{pop rdi} = \mintinline{ASM}{mov rdi, [rsp]} y \mintinline{ASM}{add rsp, 8}
	\item \mintinline{ASM}{push rdi} = \mintinline{ASM}{sub rsp ,8} y \mintinline{ASM}{mov [rsp], rdi}
\end{itemize}instrucción
Al momento de ejecutarse, la pila tiene este aspecto:


\begingroup
\setlength{\dashlinedash}{1pt}
\setlength{\dashlinegap}{2pt}
\begin{tabular}{:c:p{8cm}}
	\multicolumn{1}{c}{\ }        & [ smaller addresses, where rsp goes if you grow your own frame ] \\
	\cdashline{1-1}
	return address (back to main) & \( \leftarrow \) rsp points here                                 \\
	\cdashline{1-1}
	\makecell{caller's stack frame                                                                   \\ \dots your flag is in here \dots } &                                                                   \\
	\cdashline{1-1}
	return address (back to main) &                                                                  \\
	\cdashline{1-1}
	\dots main's frame \dots      &                                                                  \\
	\cdashline{1-1}
	\multicolumn{1}{c}{\ }        & [large addresses]                                                \\
\end{tabular}
\endgroup

\begin{flagbox}{Flag Capturada}
	pwn.college{wtu3_OlESb6TzTa1PtIXgZ9ZjUS.0FO1kDNywSMycjN4EzW}
\end{flagbox}

\subsection{Loading Stale Stack Data}\label{subsec:10-02loading_stale_stack_data}

En este caso, vamos a utilizar los bytes, que se dejarón atrás cuando una función
ya devolvió su resultado.

Cuando una función termina y regresa, los bytes que utilizo no se elimina, sino que se
quedan obsoletos, esto por un tema de eficiencia.

Cuando los datos que quedan son confidenciales, omitir ese borrado puede convertirse en una
vulnerabilidad.

\begin{flagbox}{Flag Capturada}
	pwn.college{AZKP7tCbsW41IyFS_Zzsf4hgHU2.0lN2ITNywSMycjN4EzW}
\end{flagbox}

\subsection{Stealing Stale Stack Data}\label{subsec:10-03stealing_stale_stack_data}

Se utilizará la misma idea de los datos residuales en la pila con un búfer de bytes.

En este reto, nos pide que llamemos a la función \textit{read\_flag}, y luego imprimir por
\textbf{\textit{stdout}} los bytes obsoletos de la flag provenientes del marco de la función
llamada.

\begin{flagbox}{Flag Capturada}
	pwn.college{EVdjOB2CbkZqndmuAe6SpseIdZF.01NzITNywSMycjN4EzW}
\end{flagbox}

\subsection{Reserving Your Own Frame}\label{subsec:10-04reserving_your_own_frame}

Hasta el momento, las funciones mantienen sus valores temporales en los registros;
pero una función puede necesitar más espacio de trabajo (\textit{scrathc space})
del que pueden contener los regitros.

En la arquitecutra x86, una función crea espacio de trabajo en la pila modificando
el puntero de la pila \mintinline{ASM}{rsp}, para que apunte a una dirección
más baja(\mintinline{ASM}{sub rsp, 256}).

Solo se tiene que considerar 3 cosas:

\begin{enumerate}
	\item Ahora se van a necesitar diferentes desplazamientos(\textit{offset}) relativos a \mintinline{ASM}{rsp}.
	\item Los bytes que se acaban de mover, son memoria común de la pila, y pueden contener datos dejados por código ejecutado anteriormente.
	\item Se tiene que devolver rsp a su posición original antes de retornar.
\end{enumerate}

\begin{flagbox}{Flag Capturada}
	pwn.college{kz_F2pbSUCr4BskErhXacY_APsE.0FOzITNywSMycjN4EzW}
\end{flagbox}

\subsection{Using Your Own Frame}\label{subsec:10-05using_your_own_frame}

En este desafio se buscando presentar que hagas uso de la pila que limpiaste para
luego poder hacer comparaciones, con un tamaño específico de la pila
y contar la cantidad de repetidos, buena forma depoder prácticar lo aprendiendo anteriormente.

\begin{flagbox}{Flag Capturada}
	pwn.college{UP2DMXRmvENtSlQ66pSI2CeQvbg.0FO1ATNywSMycjN4EzW}
\end{flagbox}

\subsection{Environment Variables on the Stack}\label{subsec:10-06environment_variables_on_the_stack}

Justo después de la lista de argumentos(\mintinline{C}{argc, argv}), el núcleo coloca
las variables de entorno. Al igual que \mintinline{C}{argv}, estas se almacenan en la pila
como una matriz de punteros a cadenas, en la que cada cadena incluye el valor de
la variable de la siguiente mantera \mintinline{TEXT}{HOME=/home/user/...}.

En caso el programa sea llamado sin argumentos, y una única variable de entorno llamada FLAG,
la estructura de la pila podría ser la siguiente:

\begin{verbatim}
   Address     | Contents
   +------------------------+
   | rsp + 0   | 1          | <--- argc
   +------------------------+
   | rsp + 8   | rsp + 128  |--+  argv[0]: pointer to the program name
   +------------------------+  |
   | rsp + 16  | 0          |  |  NULL (end of argv)
   +------------------------+  |
   | rsp + 24  | rsp + 200  |--+--+  envp[0]: pointer to the first env var
   +------------------------+  |  |
   | rsp + 32  | 0          |  |  |  NULL (end of envp)
   +------------------------+  |  |
                               |  |
     Address   | Contents      |  |
   +------------------------+  |  |
   | rsp + 128 | "/tmp/..." |<-+  |  the program name
   +------------------------+     |
   | ...       | ...        |     |
   +------------------------+     |
   | rsp + 200 | "FLAG=..." |<----+  the first env var: the `FLAG` variable
   +------------------------+
\end{verbatim}

Tenemos que tener en cuenta 2 cosas:

\begin{enumerate}
	\item Tanto \mintinline{C}{argv} y \mintinline{C}{envp} terminan en \mintinline{C}{NULL} al final de cada lista de punteros.
	\item Las cadenas de \mintinline{C}{envp} tiene el formato \textbf{NOMBRE=VALOR}. Por lo tanto, \mintinline{C}{envp[0]}
	      apunta a la cadena que comienza con el nombre de la primavera variable de entorno.
\end{enumerate}

\begin{flagbox}{Flag Capturada}
	flag
\end{flagbox}

\subsection{Aligning the Stack Through the Environment}\label{subsec:10-07aligning_the_stack_through_the_environment}

\begin{flagbox}{Flag Capturada}
	flag
\end{flagbox}

\subsection{Aligning the Stack Through GDB}\label{subsec:10-08aligning_the_stack_through_gdb}

\begin{flagbox}{Flag Capturada}
	flag
\end{flagbox}

\subsection{Aligning the Stack Through GDB, Generalized}\label{subsec:10-09aligning_the_stack_through_gdb_generalized}
\begin{flagbox}{Flag Capturada}
	flag
\end{flagbox}
